\section{Related work}
\label{sec:related}

\subsection{Cross-corpus speech emotion recognition}
\label{sec:related-crosscorpus}

Cross-corpus SER and strategies for handling corpus variation are established
research problems \cite{schuller2010cross}. Approaches include transfer
learning \cite{naderi2023cross}, multi-task learning \cite{fu2023cross},
training-strategy changes \cite{pastor2023cross}, adversarial adaptation
\cite{latif2023self}, and feature alignment \cite{naeeni2025feature}.
A recent \emph{Speech Communication} study combines class-sensitive feature
dictionaries with an adversarial discriminator in a semi-supervised network
\cite{zhang2025sdan}. It learns representations with additional target
supervision, unlike the frozen source-only affine maps examined here.

Frequently used corpora include IEMOCAP \cite{busso2008iemocap}, RAVDESS
\cite{livingstone2018ryerson} and CREMA-D \cite{cao2014crema}. Their label
mappings, speakers, elicitation procedures and evaluation protocols differ.
Reported scores therefore cannot be ranked without controlling those choices.
EmoBox \cite{ma2024emobox} addresses reproducible cross-corpus evaluation.
Another benchmark alone would not distinguish the present paper. Our question
concerns information available to the model-selection criterion.

\subsection{Alignment and the limits of marginal matching}
\label{sec:related-alignment}

Under covariate shift, the feature marginal changes while the conditional
label distribution is stable \cite{shimodaira2000covariate}. Domain adaptation
also studies settings outside this assumption. A standard target-error bound
contains source error, a domain divergence and the error of the best joint
hypothesis \cite{bendavid2010theory}. Its divergence is classifier-induced,
defined over the symmetric difference of the hypothesis class. Reducing a
kernel two-sample statistic alone need not reduce that divergence or the
joint-error term.

CORAL matches covariance structure \cite{sun2016coral}, and Deep CORAL makes
this a representation-learning objective \cite{sun2016deep}. MMD compares
kernel mean embeddings \cite{gretton2012kernel}; deep adaptation networks
apply multi-kernel matching to learned features \cite{long2015dan}.
Domain-adversarial networks instead use a domain classifier
\cite{ganin2016dann}. These are different adaptation objectives, not a nested
sequence of interventions on moments alone.

Marginal invariance with low source error is already known to be insufficient
for successful adaptation \cite{zhao2019invariant}. We do not claim discovery
of that failure. Our exact result concerns a validation symmetry within a
specified family of source transformations. The affine factorisation
clarifies why it does not automatically apply to the other conditions.

\subsection{Speech representations and layer choice}
\label{sec:related-ssl}

HuBERT, wav2vec 2.0 and WavLM supply the frozen features in this study
\cite{hsu2021hubert,baevski2020wav2vec,chen2022wavlm}. Final-layer pooling and
learned layer aggregation are established choices \cite{yang2021superb}.
Layer-wise probing also shows that the information represented by a speech
encoder changes with depth \cite{pasad2021layer}. The retained layer
comparisons provide context for the criterion audit. They are not presented
as the discovery that intermediate speech representations can be useful.

\subsection{Unsupervised model selection}
\label{sec:related-protocol}

Selecting hyperparameters on target test scores gives an optimistic,
non-deployable estimate. Selection bias is a general evaluation issue
\cite{cawley2010selection}. Deep Embedded Validation
\cite{you2019selection}, the study of validation practice by
\cite{ericsson2023practices}, and the reliability study of
\cite{hu2024selection} already address model selection in domain adaptation.
Their criteria use or evaluate signals beyond the source-validation score
of a single candidate. Our restricted invariance result does not invalidate
these criteria or imply that no unsupervised selector can work.

\subsection{Position of this work}
\label{sec:related-position}

The nearest methodological precedents concern adaptation model selection, not
only SER benchmarking. We contribute a concrete criterion audit: an exact
source-translation invariance, a matched-run check, and a coordinate account
that limits its extension to other affine operations. The reparameterisation
identities are elementary and the underlying selection problem is established.
Their role is to make this SER result interpretable and independently
checkable, not to claim a new general theory of adaptation.
